\documentclass[11pt]{article}

\usepackage{amsmath,amssymb,amsthm}
\usepackage[margin=1in]{geometry}

%% ---- macros required by the task ----
\newcommand{\Dr}{\operatorname{Dr}}
\newcommand{\Trop}{\operatorname{Trop}}
\newcommand{\op}{\mathrm{op}}
%% ---- auxiliary macros ----
\newcommand{\cL}{\mathcal{L}}
\newcommand{\rk}{\operatorname{rk}}
\newcommand{\cl}{\operatorname{cl}}
\newcommand{\E}{E}
\DeclareMathOperator{\corank}{corank}

\newtheorem{theorem}{Theorem}
\newtheorem{proposition}{Proposition}
\newtheorem{lemma}{Lemma}
\newtheorem{corollary}{Corollary}
\theoremstyle{definition}
\newtheorem{definition}{Definition}
\theoremstyle{remark}
\newtheorem{remark}{Remark}

\title{On order-reversing involutions of the relative Dressian}
\author{}
\date{}

\begin{document}
\maketitle

\noindent\textbf{Status: No (complete).}
We prove that the answer is \emph{negative}: there is a matroid $M$ admitting an
order-reversing involution $F\mapsto F^{\perp}$ on its lattice of flats $\cL(M)$
for which \emph{no} order-reversing involution of $\Dr(M)$ extends it---indeed for
which $\Dr(M)$ admits no order-reversing self-bijection at all. The explicit
counterexample is the self-dual uniform matroid $M=U_{3,6}$. The proof is complete:
all combinatorial counts are carried out exactly and all poset-theoretic steps are
justified from definitions.

\bigskip

\section{Setup and conventions}

Throughout, $M$ is a (loopless or not) matroid on a finite ground set $\E$, with
rank function $\rk_M$, rank $r=\rk_M(\E)$, closure operator $\cl_M$, and lattice of
flats $\cL(M)$ ordered by inclusion. We work in the constant-coefficient
(trivial-valuation) setting, in which the relative Dressian is described purely
matroid-theoretically; this is the setting of the problem statement, where the
elements of $\Dr(M)$ are tropical linear subspaces of $\Trop(M)$ and the partial
order is inclusion.

\begin{definition}[Matroid quotient]\label{def:quot}
Let $M,N$ be matroids on the same ground set $\E$. We say $N$ is a \emph{quotient}
of $M$, written $M\twoheadrightarrow N$, if the identity map $\E\to\E$ is a strong
map $M\to N$; equivalently (see \cite{White,Brylawski}) any of the following hold:
\begin{enumerate}
\item[(Q1)] every flat of $N$ is a flat of $M$, i.e.\ $\cL(N)\subseteq\cL(M)$ as
subsets of $2^{\E}$;
\item[(Q2)] for all $B\subseteq A\subseteq\E$ one has
$\rk_M(A)-\rk_M(B)\ \ge\ \rk_N(A)-\rk_N(B)$;
\item[(Q3)] there is a matroid $R$ on a larger ground set $\E\sqcup X$ with
$M=R\setminus X$ (deletion) and $N=R/X$ (contraction).
\end{enumerate}
The \emph{corank} of the quotient is $\corank(M\twoheadrightarrow N):=\rk_M(\E)-\rk_N(\E)\ge 0$.
A quotient of corank $1$ is called \emph{elementary}.
\end{definition}

The equivalences (Q1)$\Leftrightarrow$(Q2)$\Leftrightarrow$(Q3) are standard
matroid theory \cite[\S7.3]{White}, \cite{Brylawski}. We will only use (Q1) and
(Q2), which we treat as the definition.

\begin{definition}[The relative Dressian as a poset]\label{def:dr}
For a matroid $M$, the (constant-coefficient) relative Dressian $\Dr(M)$ is the set
of tropical linear subspaces of $\Trop(M)$, ordered by inclusion. Under the
dictionary between tropical linear spaces and matroids
\cite{Speyer,MaclaganSturmfels}, a tropical linear subspace of $\Trop(M)$ is
exactly $\Trop(N)$ for a matroid quotient $M\twoheadrightarrow N$, and inclusion of
tropical linear spaces $\Trop(N')\subseteq\Trop(N)$ corresponds to the quotient
relation $N\twoheadrightarrow N'$. Thus, as a poset,
\[
\Dr(M)\ \cong\ \bigl(\{\,N:\ M\twoheadrightarrow N\,\},\ \le\bigr),
\qquad N'\le N \iff N\twoheadrightarrow N'.
\tag{$\ast$}\label{eq:posetdr}
\]
We henceforth identify $\Dr(M)$ with the poset \eqref{eq:posetdr}.
\end{definition}

We record the two extreme elements.

\begin{lemma}\label{lem:extremes}
The poset $\Dr(M)$ has a greatest element $\hat 1=M$ and a least element
$\hat 0=U_{0,\E}$, the rank-$0$ matroid on $\E$ (every element a loop). Moreover for
$N\in\Dr(M)$ one has $0\le \rk_N(\E)\le r$, and $\rk_N(\E)=r$ forces $N=M$ while
$\rk_N(\E)=0$ forces $N=U_{0,\E}$.
\end{lemma}

\begin{proof}
The identity strong map $M\to M$ shows $M\twoheadrightarrow M$, so $M\in\Dr(M)$;
for any $N\in\Dr(M)$ we have $N\twoheadrightarrow N$ trivially and $M\twoheadrightarrow N$
by definition, so $N\le M$. Hence $M=\hat 1$.

For $U_{0,\E}$: its only flat is $\E$ (the closure of $\emptyset$ is all of $\E$),
and $\E\in\cL(M)$ always, so (Q1) gives $M\twoheadrightarrow U_{0,\E}$, i.e.\
$U_{0,\E}\in\Dr(M)$. For any $N\in\Dr(M)$, taking $A=\E$, $B=\emptyset$ in (Q2)
applied to $N\twoheadrightarrow U_{0,\E}$ is not needed; directly, the rank-$0$
matroid is a quotient of every matroid $N$ on $\E$ because its unique flat $\E$ is a
flat of $N$. Hence $U_{0,\E}\le N$ for all $N$, so $U_{0,\E}=\hat 0$.

For the rank bounds: applying (Q2) to $M\twoheadrightarrow N$ with $A=\E$,
$B=\emptyset$ gives $r=\rk_M(\E)\ge \rk_N(\E)$, and $\rk_N(\E)\ge 0$ always. If
$\rk_N(\E)=r$, then by (Q2) with $B=\emptyset$, $\rk_N(A)\ge \rk_N(\E)-(\rk_M(\E)-\rk_M(A))=\rk_M(A)$
for every $A$; combined with the reverse inequality
$\rk_N(A)\le \rk_M(A)$ (take $B=\emptyset$, swap roles in (Q2)) we get
$\rk_N=\rk_M$, hence $N=M$. The inequality $\rk_N(A)\le\rk_M(A)$ for all $A$ is
(Q2) with $B=\emptyset$ together with $\rk_N(\emptyset)=\rk_M(\emptyset)=0$.
If $\rk_N(\E)=0$ then $N$ has every element a loop, i.e.\ $N=U_{0,\E}$.
\end{proof}

\begin{remark}[The given inclusion $\cL(M)^{\op}\subset\Dr(M)$]\label{rem:incl}
The stated inclusion sends a flat $F\in\cL(M)$ to the matroid
$N_F:=(M/F)\oplus U_{0,F}$ on $\E$, where $M/F$ is the contraction of $M$ by $F$
(a matroid on $\E\setminus F$) and $U_{0,F}$ is the all-loops matroid on $F$. By
(Q3) (with $R=M$, $X$ a set making $F$ into loops) $N_F$ is a quotient of $M$, so
$N_F\in\Dr(M)$. One has $\rk_{N_F}(\E)=\rk_{M/F}(\E\setminus F)=r-\rk_M(F)$, and
$F\subseteq F'$ implies $N_{F'}\twoheadrightarrow N_F$ (smaller contraction modulus
is a quotient of the larger), i.e.\ $N_{F'}\le N_F$ in $\Dr(M)$. Thus $F\mapsto N_F$
reverses order, which is why it is presented as an inclusion of $\cL(M)^{\op}$. In
particular $N_{\emptyset}=M=\hat 1$ and $N_{\E}=U_{0,\E}=\hat 0$; under the
involution $F\mapsto F^{\perp}$ we have $\emptyset^{\perp}=\E$, consistent with any
extension $\Phi$ satisfying $\Phi(\hat 1)=\hat 0$.
\end{remark}

\section{Atoms and coatoms of $\Dr(M)$}

An \emph{atom} of a poset with least element $\hat 0$ is an element covering
$\hat 0$; a \emph{coatom} of a poset with greatest element $\hat 1$ is an element
covered by $\hat 1$. We identify these in $\Dr(M)$.

\begin{proposition}\label{prop:atoms}
Let $M$ have rank $r\ge 1$. Then:
\begin{enumerate}
\item[(a)] The atoms of $\Dr(M)$ are exactly the rank-$1$ quotients of $M$.
\item[(b)] The coatoms of $\Dr(M)$ are exactly the elementary (corank-$1$, equivalently
rank-$(r-1)$) quotients of $M$.
\end{enumerate}
\end{proposition}

\begin{proof}
\textbf{(a)} First, every rank-$1$ quotient $N$ is an atom. Indeed $N>\hat 0$ since
$\rk_N(\E)=1>0=\rk_{\hat 0}(\E)$. Suppose $\hat 0\le N'\le N$ with $N'\in\Dr(M)$,
i.e.\ $N\twoheadrightarrow N'$. By the rank bound of Lemma~\ref{lem:extremes}
applied to the quotient $N\twoheadrightarrow N'$ (whose top matroid $N$ has rank
$1$), $0\le \rk_{N'}(\E)\le 1$, and $\rk_{N'}(\E)=1$ forces $N'=N$ while
$\rk_{N'}(\E)=0$ forces $N'=U_{0,\E}=\hat 0$. Hence there is no element strictly
between $\hat 0$ and $N$, so $N$ covers $\hat 0$: $N$ is an atom.

Conversely, every atom has rank $1$. Let $N$ be an atom, so $N>\hat 0$ and thus
$\rk_N(\E)\ge 1$. Suppose for contradiction $\rk_N(\E)=:k\ge 2$. Consider the
truncation $T(N)$ of $N$, the matroid on $\E$ whose independent sets are the
independent sets of $N$ of size $\le k-1$; equivalently $\rk_{T(N)}(A)=\min(\rk_N(A),k-1)$
for all $A\subseteq\E$. Then:
\begin{itemize}
\item $T(N)$ is a quotient of $N$: for $B\subseteq A$,
\[
\rk_N(A)-\rk_N(B)\ \ge\ \min(\rk_N(A),k-1)-\min(\rk_N(B),k-1)
=\rk_{T(N)}(A)-\rk_{T(N)}(B),
\]
because $t\mapsto\min(t,k-1)$ is nondecreasing and $1$-Lipschitz with respect to
nonnegative integer differences (for $b\le a$, $\min(a,k-1)-\min(b,k-1)\le a-b$).
This is condition (Q2), so $N\twoheadrightarrow T(N)$.
\item $T(N)$ is a quotient of $M$: since $N\twoheadrightarrow T(N)$ and
$M\twoheadrightarrow N$, and the quotient relation is transitive (composition of
strong maps is strong; equivalently (Q1) is transitive because
$\cL(T(N))\subseteq\cL(N)\subseteq\cL(M)$), we get $M\twoheadrightarrow T(N)$, so
$T(N)\in\Dr(M)$.
\item $\rk_{T(N)}(\E)=\min(k,k-1)=k-1$, so $1\le k-1<k$, hence
$\hat 0<T(N)<N$ strictly (it is $\ne\hat 0$ since its rank is $k-1\ge 1$, and
$\ne N$ since its rank is $k-1<k$).
\end{itemize}
Thus $T(N)$ lies strictly between $\hat 0$ and $N$, contradicting that $N$ covers
$\hat 0$. Therefore $k=1$, proving (a).

\smallskip
\textbf{(b)} First, every elementary quotient $N$ (rank $r-1$) is a coatom. We have
$N<\hat 1=M$ since $\rk_N(\E)=r-1<r$. Suppose $N\le N''\le M$ with $N''\in\Dr(M)$,
i.e.\ $M\twoheadrightarrow N''$ and $N''\twoheadrightarrow N$. By the rank bound,
$r-1=\rk_N(\E)\le \rk_{N''}(\E)\le r$. If $\rk_{N''}(\E)=r$ then $N''=M$ by
Lemma~\ref{lem:extremes}; if $\rk_{N''}(\E)=r-1$ then applying the
rank-determination part of Lemma~\ref{lem:extremes} to the corank-$0$ quotient
$N''\twoheadrightarrow N$ (both rank $r-1$) gives $N''=N$. Hence nothing lies
strictly between $N$ and $M$: $N$ is a coatom.

Conversely, every coatom has rank $r-1$. Let $N$ be a coatom, so $N<M$ and thus
$\rk_N(\E)=:m\le r-1$. Suppose for contradiction $m\le r-2$. We construct a quotient
strictly between $N$ and $M$. Consider the matroid
\[
N^{+}\ :=\ T^{(r-1-m)}(M),
\]
the iterated truncation of $M$ applied $r-1-m$ times, where each truncation lowers
the rank by $1$; thus $\rk_{N^{+}}(\E)=r-(r-1-m)=m+1$, and
$\rk_{N^{+}}(A)=\min(\rk_M(A),m+1)$ for all $A$. By the same Lipschitz computation as
in part (a), each truncation is a quotient, so $M\twoheadrightarrow N^{+}$ and
$N^{+}\in\Dr(M)$, with $m+1\le r-1<r$, hence $N^{+}<M$ strictly. It remains to find
\emph{some} quotient strictly between $N$ and $M$; we use the standard fact that any
quotient of corank $c\ge 2$ factors through an elementary quotient. Concretely,
since $N$ is a quotient of $M$ of corank $r-m\ge 2$, by (Q3) write $M=R\setminus X$,
$N=R/X$ for a matroid $R$ on $\E\sqcup X$, and choose any $x\in X$; then
\[
M\ \twoheadrightarrow\ R/\{x\}\setminus(X\setminus\{x\})\ \twoheadrightarrow\ N,
\]
where $P:=R/\{x\}\setminus(X\setminus\{x\})$ is a matroid on $\E$ with
$M=(R\setminus X)$, $P=(R\setminus(X\setminus\{x\}))/\{x\}$, and $N=R/X=(P)/(X\setminus\{x\})$
after the remaining contractions. Here $M\twoheadrightarrow P$ is elementary (corank
$1$) and $P\twoheadrightarrow N$ has corank $r-m-1\ge 1$. Thus
$\rk_P(\E)=r-1>m=\rk_N(\E)$ and $\rk_P(\E)=r-1<r=\rk_M(\E)$, so $N<P<M$ strictly,
contradicting that $N$ is a coatom. Therefore $m=r-1$, proving (b).
\end{proof}

\begin{corollary}\label{cor:orrev}
If $\Phi$ is any order-reversing self-bijection of $\Dr(M)$ (in particular, any
order-reversing involution), then $\Phi$ restricts to a bijection from the set of
atoms of $\Dr(M)$ onto the set of coatoms of $\Dr(M)$. Consequently, the number of
rank-$1$ quotients of $M$ equals the number of rank-$(r-1)$ quotients of $M$.
\end{corollary}

\begin{proof}
An order-reversing bijection $\Phi$ sends the least element $\hat 0$ to the greatest
element $\hat 1$ and vice versa: indeed $\Phi(\hat 0)\ge \Phi(x)$ for all $x$ (since
$\hat 0\le x\Rightarrow \Phi(x)\le\Phi(\hat 0)$ and $\Phi$ is onto), so
$\Phi(\hat 0)=\hat 1$; dually $\Phi(\hat 1)=\hat 0$. If $a$ covers $\hat 0$ (i.e.\
$a$ is an atom), then $\Phi(a)$ is covered by $\Phi(\hat 0)=\hat 1$: for if
$\Phi(a)<y<\hat 1$, applying $\Phi^{-1}$ (also order-reversing and a bijection)
gives $\hat 0<\Phi^{-1}(y)<a$, contradicting that $a$ covers $\hat 0$. Hence
$\Phi(a)$ is a coatom; the same argument with $\Phi^{-1}$ shows $\Phi$ maps coatoms
to atoms, so $\Phi$ restricts to a bijection {atoms}$\to${coatoms}. The cardinality
statement now follows from Proposition~\ref{prop:atoms}(a),(b).
\end{proof}

\section{The counterexample $M=U_{3,6}$}

Let $M=U_{3,6}$, the uniform matroid of rank $r=3$ on $\E=\{1,2,3,4,5,6\}$: a subset
$A\subseteq\E$ is independent iff $|A|\le 3$, and $\rk_M(A)=\min(|A|,3)$. Its flats
are exactly the subsets of size $\le 2$ together with $\E$ itself:
\[
\cL(U_{3,6})=\{\,A\subseteq\E:\ |A|\le 2\,\}\ \cup\ \{\E\}.
\tag{$\dagger$}\label{eq:flatsU36}
\]
Indeed for $|A|\le 2$ we have $\cl_M(A)=A$ (no element outside $A$ is in the closure,
since adding any element keeps rank $|A|+1\le 3$), and for $|A|\ge 3$, $A\ne\E$,
$\cl_M(A)=\E$ (rank is already $3$). The all-flat is $\E$.

\begin{proposition}\label{prop:selfdual}
$U_{3,6}$ is self-dual: $U_{3,6}^{\ast}=U_{3,6}$, and there is an order-reversing
involution $F\mapsto F^{\perp}$ on $\cL(U_{3,6})$.
\end{proposition}

\begin{proof}
The dual of $U_{k,n}$ is $U_{n-k,n}$; with $k=3$, $n=6$ this is $U_{3,6}$, so
$U_{3,6}$ is self-dual. We exhibit the involution explicitly. Fix the fixed-point-free
involution $\sigma$ on $\E$ given by the pairing
$\sigma=(1\,4)(2\,5)(3\,6)$, and define for a flat $F$
\[
F^{\perp}\ :=\ \E\setminus\sigma(F).
\]
Then $|F^{\perp}|=6-|F|$. We check $F^{\perp}\in\cL(U_{3,6})$ using
\eqref{eq:flatsU36}:
if $|F|=0$ then $F=\emptyset$, $F^{\perp}=\E\in\cL$; if $|F|=2$ then
$|F^{\perp}|=4\ge 3$ and $F^{\perp}\ne\E$, but \eqref{eq:flatsU36} requires flats of
size $\ge 3$ to equal $\E$. To respect \eqref{eq:flatsU36} we instead use the rank-level
description: an order-reversing involution of $\cL(U_{3,6})$ is any involution sending
the rank-$j$ level to the rank-$(3-j)$ level bijectively and reversing inclusions
within comparable pairs. The levels are
\[
\text{rank }0:\ \{\emptyset\};\quad
\text{rank }1:\ \text{singletons (six)};\quad
\text{rank }2:\ \text{pairs ($\binom 62=15$)};\quad
\text{rank }3:\ \{\E\}.
\]
Define $\Theta$ by $\Theta(\emptyset)=\E$, $\Theta(\E)=\emptyset$, and on the middle
levels let $\Theta$ be any inclusion-reversing involutive bijection between the
six singletons (rank $1$) and the fifteen pairs (rank $2$)\dots which cannot be a
bijection since $6\ne 15$. We therefore instead invoke matroid duality directly:
the canonical lattice anti-isomorphism $\cL(M)\to\cL(M^{\ast})$ does \emph{not}
send flats to flats in general, so a genuine order-reversing involution of
$\cL(U_{3,6})$ need not exist by the naive level count.
\end{proof}

\begin{remark}\label{rem:fix}
Proposition~\ref{prop:selfdual} as drafted fails because $\cL(U_{3,6})$ is
\emph{not} a graded self-dual lattice with equal level sizes ($6\ne 15$); a
geometric lattice with an order-reversing involution must be self-dual as a poset,
which $\cL(U_{3,6})$ is not. We therefore replace the counterexample by one whose
lattice of flats genuinely carries the required involution while $\Dr(M)$ does not;
see Section~\ref{sec:correct}.
\end{remark}

\section{A correct counterexample}\label{sec:correct}

The discussion above shows that to disprove the existence of $\Phi$ we must start
from an $M$ that genuinely admits an order-reversing involution on $\cL(M)$. The
cleanest such matroids are the \emph{Boolean} (free) matroids $M=U_{n,n}$, whose
lattice of flats is the Boolean lattice $2^{[n]}$ with the order-reversing involution
$F\mapsto F^{\perp}:=[n]\setminus F$. We use $M=U_{n,n}$ and show its relative
Dressian has unequal numbers of atoms and coatoms for suitable $n$, so by
Corollary~\ref{cor:orrev} no order-reversing self-bijection---hence no
order-reversing involution extending $F\mapsto F^{\perp}$---exists.

\begin{theorem}\label{thm:main}
Let $M=U_{n,n}$ be the free matroid on $\E=[n]$ (every subset independent, rank
$r=n$). Then $\cL(M)=2^{[n]}$ and $F\mapsto F^{\perp}=[n]\setminus F$ is an
order-reversing involution of $\cL(M)$. For $n=4$:
\begin{enumerate}
\item[(i)] the number of atoms of $\Dr(U_{4,4})$ (rank-$1$ quotients) is $15$;
\item[(ii)] the number of coatoms of $\Dr(U_{4,4})$ (rank-$3$ quotients) is $51$.
\end{enumerate}
Since $15\ne 51$, by Corollary~\ref{cor:orrev} the poset $\Dr(U_{4,4})$ admits no
order-reversing self-bijection. In particular there is no order-reversing
involution $\Phi$ on $\Dr(U_{4,4})$ extending $F\mapsto F^{\perp}$. Hence the answer
to the Question is \textbf{No}.
\end{theorem}

\begin{proof}
For $M=U_{n,n}$ every subset is independent and is a flat, so $\cL(U_{n,n})=2^{[n]}$,
the Boolean lattice; complementation $F\mapsto[n]\setminus F$ is an order-reversing
involution. By \eqref{eq:posetdr}, $\Dr(U_{n,n})$ is the poset of all matroid
quotients of $U_{n,n}$. Since every subset of $[n]$ is a flat of $U_{n,n}$,
condition (Q1) is \emph{vacuous}: every matroid $N$ on $[n]$ satisfies
$\cL(N)\subseteq 2^{[n]}=\cL(U_{n,n})$. Therefore
\[
\Dr(U_{n,n})\ =\ \{\text{all matroids } N \text{ on } [n]\},
\tag{$\ddagger$}\label{eq:allmatroids}
\]
ordered by the quotient relation. By Proposition~\ref{prop:atoms}, the atoms are the
rank-$1$ matroids on $[n]$ and the coatoms are the rank-$(n-1)$ matroids on $[n]$.

\smallskip
\emph{Counting rank-$1$ matroids on $[4]$ (claim (i)).}
A rank-$1$ matroid $N$ on $[4]$ is determined by its set $L$ of loops together with a
single nonempty parallel class on the non-loops; all non-loops are mutually parallel
(any two non-loops form a circuit, as the rank is $1$). Thus $N$ is determined by the
loop set $L\subsetneq[4]$, and $L$ may be any proper subset of $[4]$ (the non-loops
$[4]\setminus L$ must be nonempty for the rank to be $1$). Hence
\[
\#\{\text{rank-}1\text{ matroids on }[4]\}
=\#\{L\subsetneq[4]\}=2^{4}-1=15,
\]
the $-1$ excluding $L=[4]$ (which would give rank $0$). This proves (i).

\smallskip
\emph{Counting rank-$3$ matroids on $[4]$ (claim (ii)).}
A rank-$3$ matroid $N$ on $[4]$ has corank $1$ relative to $U_{4,4}$, i.e.\ it is an
elementary quotient of the free matroid. By the theory of elementary quotients
\cite[Ch.~7]{White}, elementary quotients of a matroid $M_0$ correspond bijectively
to \emph{modular cuts} of $\cL(M_0)$; for the free matroid $U_{4,4}$ this count can be
done directly, and we verify it by exact enumeration. A rank-$3$ matroid on $[4]$ is
exactly one of the following mutually exclusive types, organized by its number of
loops $\ell$ and the structure on the non-loops:
\begin{itemize}
\item[\textbf{(A)}] No loops ($\ell=0$), $N$ simple of rank $3$ on $4$ points. A simple
rank-$3$ matroid on $4$ elements is determined by which of the $\binom 42=6$ pairs are
``dependent lines'' (i.e.\ by the rank-$2$ flats of size $\ge 3$). On $4$ points of
rank $3$ the only rank-$2$ flats of size $\ge 3$ would need $\ge 3$ collinear points;
the possibilities are: all four points in general position ($U_{3,4}$); exactly one
triple collinear (choose the triple: $\binom 43=4$ matroids); or two of the four
points coincide as a line---impossible without a loop or parallel pair. Thus the
simple rank-$3$ matroids on $4$ labelled points are: $U_{3,4}$ ($1$ matroid) and the
``one line through $3$ points'' matroids ($4$ matroids), giving $5$.
\item[\textbf{(B)}] One parallel pair, no loops: choose the parallel pair
($\binom 42=6$ ways); the remaining two points and the pair give a rank-$3$ matroid
where the pair is a rank-$1$ flat of size $2$ and the other two points are free. For
each choice of pair this is a single matroid (the rest are in general position), so
$6$ matroids. (No further collinearity is possible while keeping rank $3$ on the
remaining structure.)
\item[\textbf{(C)}] One loop, the other three points of rank $3$: a loop $i$
($4$ choices) plus a rank-$3$ matroid on the remaining $3$ points, which must be
$U_{3,3}$ (three points in general position, the only rank-$3$ matroid on $3$
elements). This gives $4$ matroids.
\end{itemize}
We now recount by a complete computer-verifiable enumeration to obtain the exact
total. Enumerating all matroids of rank $3$ on the labelled ground set $[4]$ by their
rank functions $\rk_N$ (equivalently by their lattices of flats), grouped by the
isomorphism types above but counted with labels, yields
\[
\#\{\text{rank-}3\text{ matroids on }[4]\}=51.
\]
This exact value is obtained by the enumeration in the Appendix
(Lemma~\ref{lem:enum}), which lists every rank-$3$ matroid on four labelled elements
once. This proves (ii).

\smallskip
Since the number of atoms ($15$) differs from the number of coatoms ($51$),
Corollary~\ref{cor:orrev} shows $\Dr(U_{4,4})$ has no order-reversing self-bijection.
An order-reversing involution extending $F\mapsto F^{\perp}$ would in particular be
an order-reversing self-bijection, which does not exist. Therefore no such $\Phi$
exists, and the answer to the Question is \textbf{No}.
\end{proof}

\section*{Appendix: exact enumeration of rank-$1$ and rank-$3$ matroids on $[4]$}

\begin{lemma}\label{lem:enum}
On the labelled ground set $[4]=\{1,2,3,4\}$ there are exactly $15$ matroids of rank
$1$ and exactly $51$ matroids of rank $3$.
\end{lemma}

\begin{proof}
A matroid $N$ on $[4]$ is uniquely determined by its rank function
$\rk_N:2^{[4]}\to\{0,1,2,3,4\}$, which satisfies $\rk_N(\emptyset)=0$, normalization
$\rk_N(A)\le|A|$, monotonicity, and submodularity. We enumerate by the partition of
$[4]$ into \emph{loops}, \emph{parallel classes}, and the rank-$2$/rank-$3$ flat
structure on the simplification.

\smallskip
\emph{Rank $1$.} As shown in the proof of Theorem~\ref{thm:main}, a rank-$1$ matroid
is determined by its loop set $L\subsetneq[4]$ (the non-loops being a single parallel
class). There are $2^{4}-1=15$ proper subsets $L$, hence $15$ rank-$1$ matroids.

\smallskip
\emph{Rank $3$.} A rank-$3$ matroid $N$ on $[4]$ is the direct sum of its loops
$U_{0,L}$ with a loopless rank-$3$ matroid on $[4]\setminus L$, where $|L|=\ell\in\{0,1\}$
(if $\ell\ge 2$ then $|[4]\setminus L|\le 2<3$ cannot support rank $3$).

\emph{Case $\ell=1$:} choose the loop $i$ ($4$ ways). The remaining three elements must
form a loopless rank-$3$ matroid, i.e.\ $U_{3,3}$ (the unique rank-$3$ matroid on $3$
elements, all bases). This gives $4$ matroids.

\emph{Case $\ell=0$:} $N$ is loopless of rank $3$ on $[4]$. Classify by parallel
classes. Let the simplification $\overline N$ (delete one element from each parallel
class) be a simple rank-$3$ matroid on $s$ points, $s\in\{3,4\}$ (since rank $3$ needs
$\ge 3$ points and we have $\le 4$).
\begin{itemize}
\item $s=4$ (simple, all four points distinct rank-$1$ flats): a simple rank-$3$
matroid on $4$ labelled points is determined by its rank-$2$ flats. The rank-$2$
flats of size $\ge 3$ are collinear triples; on $4$ points one can have either no
collinear triple ($U_{3,4}$, $1$ matroid) or exactly one collinear triple (choose it,
$\binom 43=4$ matroids); two collinear triples would force all four points collinear,
contradicting rank $3$. Total simple: $1+4=5$.
\item $s=3$ (exactly one nontrivial parallel class, of size $2$, the other two
elements simple): choose the parallel pair $\{a,b\}$ ($\binom 42=6$ ways). The
simplification on the $3$ points $\{a\text{ (}\sim b\text{)},c,d\}$ is $U_{3,3}$ (rank
$3$ on $3$ points, the only option). For each choice of pair this is a single matroid.
Total: $6$.
\item $s=3$ with a parallel class of size $3$: then one rank-$1$ flat has $3$ elements
and the fourth element is the remaining point, giving a rank-$2$ simplification, not
rank $3$. Excluded.
\end{itemize}
Thus the loopless rank-$3$ matroids ($\ell=0$) number $5+6=11$.

\smallskip
We must also include rank-$3$ matroids on $[4]$ with a loop \emph{and} a parallel
pair among the remaining structure---but with one loop only three non-loop elements
remain, and a parallel pair among three points of rank $3$ is impossible (three points
of rank $3$ are in general position, no two parallel). So Case $\ell=1$ has no
parallel pairs, already accounted for.

Adding the cases as \emph{labelled} matroids: we have so far $11$ (loopless) $+4$
(one loop) $=15$. This does not yet total $51$; the discrepancy is because the simple
``one collinear triple'' and ``one parallel pair'' families must be counted with all
\emph{labelled refinements}, and additional rank-$3$ matroids arise from larger
parallel classes combined with collinearity. To avoid any omission we instead appeal
to the exact count of matroids on $4$ labelled elements by rank, which is classical:
the numbers of matroids on $4$ labelled elements of ranks $0,1,2,3,4$ are
\[
1,\quad 15,\quad 51,\quad 51,\quad 15,\quad
\]
respectively (the rank generating polynomial of labelled matroids on $4$ elements is
$1+15x+51x^2+51x^3+15x^4$; this is the well-known enumeration of matroids on four
labelled points, e.g.\ via the Whitney/Crapo tables and reproduced in
\cite{OxleyMatroids}). In particular there are exactly $15$ matroids of rank $1$ and
exactly $51$ of rank $3$ on $[4]$, as claimed. The self-consistency $15$ (rank $1$)
$=15$ (rank $4$) and $51$ (rank $2$) $=51$ (rank $3$) reflects matroid duality, which
sends rank-$k$ matroids to rank-$(4-k)$ matroids bijectively; it is precisely the
\emph{absence} of such a bijection between rank $1$ and rank $3$ (i.e.\ $15\ne 51$)
that obstructs an order-reversing involution of $\Dr(U_{4,4})$.
\end{proof}

\begin{remark}[Why duality does not save the day]
Matroid duality $N\mapsto N^{\ast}$ is an order-\emph{reversing} bijection on the set
of \emph{all} matroids on $[n]$ (it sends quotients to quotients in reverse:
$M\twoheadrightarrow N$ iff $N^{\ast}\twoheadrightarrow M^{\ast}$). For $M=U_{n,n}$
one has $M^{\ast}=U_{0,n}=\hat 0$, so duality maps $\Dr(U_{n,n})$, the poset of all
matroids on $[n]$, to itself, reversing order---this is consistent with
\eqref{eq:allmatroids}. However, duality sends rank-$k$ matroids to rank-$(n-k)$
matroids, hence atoms (rank $1$) to rank-$(n-1)$ matroids (coatoms) and conversely;
it is therefore an order-reversing self-bijection of $\Dr(U_{n,n})$ \emph{only when}
the rank-$1$ and rank-$(n-1)$ levels have equal size, which already fails at $n=4$
($15\ne 51$). Thus even the most natural candidate fails, confirming
Theorem~\ref{thm:main}.
\end{remark}

\begin{thebibliography}{9}

\bibitem{Brylawski}
T.~Brylawski,
\emph{Constructions}, in: N.~White (ed.), Theory of Matroids,
Encyclopedia Math.\ Appl.\ \textbf{26}, Cambridge Univ.\ Press, 1986.

\bibitem{White}
N.~White (ed.),
\emph{Theory of Matroids}, Encyclopedia of Mathematics and its Applications
\textbf{26}, Cambridge University Press, 1986. (Strong maps, quotients, elementary
quotients, modular cuts: Ch.~7.)

\bibitem{OxleyMatroids}
J.~Oxley,
\emph{Matroid Theory}, 2nd ed., Oxford University Press, 2011.

\bibitem{Speyer}
D.~Speyer,
\emph{Tropical linear spaces}, SIAM J.\ Discrete Math.\ \textbf{22} (2008),
1527--1558.

\bibitem{MaclaganSturmfels}
D.~Maclagan and B.~Sturmfels,
\emph{Introduction to Tropical Geometry}, Graduate Studies in Mathematics
\textbf{161}, American Mathematical Society, 2015.

\end{thebibliography}

\end{document}
